\documentclass{article}
\usepackage[utf8]{inputenc}
\usepackage{xcolor}
\usepackage{hyperref}

\usepackage{rotating}

\usepackage{tabls}
\usepackage{booktabs}
\usepackage{amsmath}
\usepackage{amssymb}
\usepackage{amsthm}
\usepackage{amsfonts}
\usepackage{multicol}
\usepackage{enumitem}
\usepackage{microtype}
\usepackage{tikz}
\usepackage{pgfplots}
\usetikzlibrary{positioning}
\usepackage{multirow}

\usepackage{comment}

\usepackage{graphicx}
\usepackage{wrapfig}

\theoremstyle{plain}
\newtheorem{theorem}{Theorem}
\newtheorem*{theorem*}{Theorem}
\newtheorem{lemma}[theorem]{Lemma}
\newtheorem*{lemma*}{Lemma}
\newtheorem{proposition}[theorem]{Proposition}
\newtheorem{corollary}[theorem]{Corollary}

\theoremstyle{box}
\newtheorem{definition}[theorem]{Definition}
\newtheorem*{axiom*}{Axiom}
\newtheorem{dfn}[theorem]{Definition}
\newtheorem*{definition*}{Definition}
\newtheorem*{exercise*}{Exercise}
\newtheorem{observation}[theorem]{Observation}
\newtheorem{remark}[theorem]{Remark}
\newtheorem{example}[theorem]{Example}
\newtheorem{question}[theorem]{Question}
\newtheorem*{prep-problems}{Preparation Problems}

\newcommand{\pd}[3]{\frac{\partial^{#3} {#1}}{\partial {#2}^{#3}}}
\newcommand{\fpd}[2]{\frac{\partial {#1}}{\partial {#2}}}
\newcommand{\diff}[2]{\frac{d {#1}}{d {#2}}}


\begin{document}
\section{Objectives}
\begin{itemize}
\item Deepen and broaden our understanding of functions.
	\begin{itemize}
	\item Identify the domain and range of a probability mass function.
	\item Identify the domain and range of some common probability and statistics functions.
	\item Identify the parameters of some common probability and statistics functions.
	\item Explore some common functions from probability and statistics using R.
	\end{itemize}
\item Write down the pmf for some discrete random variables.
\item Use the pmf to calculate probability.
\item Recognize that notation helps us communicate our mathematical thinking about functions.
	\begin{itemize}
	\item Use expanded notation and product notation to write a factorial.
	\end{itemize}
\item 
\end{itemize}
\subsection{Resources:}
\subsection{Announcements:}


%\subsection{table of contents}

\newpage
\section{Agenda}
\subsection{Graphing a PMF}
Consider the experiment where two 6-sided dice are rolled and we sum the dice. Let $S$ be the random variable that represents the value of the sum of the dice. The pmf of $S$ is
      $$f(s) = 
    \begin{cases}
    \frac{1}{36} & s=2,12 \\
    \frac{1}{18} & s=3,11 \\
    \frac{1}{12} & s=4,10 \\
    \frac{1}{9} & s=5,9 \\
    \frac{5}{36} & s=6,8 \\
    \frac{1}{6} & s=7 \\
    0 & \text{otherwise}.
    \end{cases}$$

Follow up questions from Prep:
\begin{itemize}
\item What is the sample space for the 2 dice experiment?
\item What is the domain of $f$?
\item What is the range of $f$?
\item What does $f$ look like?\\
Let's graph $f$ (both by hand and in R).
\end{itemize}
 

 



\newpage
\subsection{Common Functions from Probability and Statistics}
\begin{definition*}[Factorial]
The product of positive integers from 1 to $n$ is called $n$ factorial and is written $n$!
\end{definition*}

\noindent Write 3! in expanded notation and product notation.\\
Write n! in expanded notation and product notation.\\
Talk about the factorial function in R.

\begin{exercise*}[]
For each of the functions, do the following:
\begin{itemize}
\item Identify the domain.
\item Identify the parameters.
\item Plot the function in R.
\item Identify the range.
\item Explore what happens when you change the parameters.
\item Describe what changing each of the parameters does to the graph of $f$.
\end{itemize}
\vfill

Below are a some common functions seen in Probability and Statistics:\\

$\displaystyle{
f_1(x) = \left\{
        \begin{array}{ll}
            \frac{n!}{x!(n-x)!}p^x(1-p)^{(n-x)} & \quad x = 0, 1, 2, 3, ... , n \\
            0 & \quad \text{otherwise}
        \end{array}
    \right.}$\\
\vfill

$\displaystyle{
f_2(x) = \left\{
        \begin{array}{ll}
            e^{-\lambda}\frac{\lambda^x}{x!} & \quad x = 0, 1, 2, 3, ... , n \\
            0 & \quad \text{otherwise}
        \end{array}
    \right.}$
\vfill

%$\displaystyle{
%f(x) = \left\{
%        \begin{array}{ll}
%            \frac{1}{b-a} & \quad a < x < b \\
%            0 & \quad \text{otherwise}
%        \end{array}
%    \right.}$
%\vfill

$\displaystyle{
f_3(x) = \left\{
        \begin{array}{ll}
            \lambda e^{-\lambda x} & \quad x > 0 \\
            0 & \quad x \leq 0
        \end{array}
    \right.}$
\vfill

$\displaystyle{
f_4(x) = \frac{1}{\sqrt{2\pi\sigma^2}}e^{-\frac{(x-\mu)^2}{2\sigma^2}}}$
\end{exercise*}



\newpage
\subsection{PREP for Monday}
Read this data into R and plot scatterplots of x1 vs. out1, x1 vs. out2, and x1 vs. out3.
\begin{exercise*}[]
\end{exercise*}



 
\end{document}